\documentclass[12pt]{article}
\usepackage[margin=1in]{geometry}
\usepackage[all]{xy}


\usepackage{amsmath,amsthm,amssymb,color,latexsym}
\usepackage{geometry}        
\geometry{letterpaper}  
\usepackage{boxedminipage}
\usepackage{graphicx}
\usepackage{hyperref}
\hypersetup{
    colorlinks=true,
    citecolor=blue,
    linkcolor=blue,
    filecolor=cyan,
    urlcolor=magenta,
}
\usepackage[capitalize,nameinlink]{cleveref}

\newtheorem{problem}{Problem}

\newenvironment{solution}[1][\it{Solution}]{\textbf{#1. } }{$\square$}
\newcommand{\pointsbox}[1]{%
  \colorbox{red!70}{\strut\textcolor{white}{\sffamily\bfseries\ \ #1\ \ }}%
}
\newcommand{\pagetag}[1]{%
  \colorbox{red!70}{\strut\textcolor{white}{\sffamily\bfseries\ \ #1\ \ }}%
}
\newcommand{\striptitle}[2]{%
  \noindent
  \begin{minipage}{0.07\linewidth}\centering \pagetag{#1}\end{minipage}%
  \hfill
  \begin{minipage}{0.89\linewidth}
    {\large\scshape #2}
  \end{minipage}\par\vspace{0.75em}
}
% --- Environments ---
\theoremstyle{plain}
\newtheorem{theorem}{Theorem}
\theoremstyle{definition}
\newtheorem{definition}[theorem]{Definition}
\newtheorem{remark}[theorem]{Remark}

\newenvironment{boxedalgo}
  {\begin{center}\begin{boxedminipage}{1\linewidth}}
  {\end{boxedminipage}\end{center}}

\newenvironment{inlinealgo}[1]
  {\noindent\hrulefill\\\noindent{\bf {#1}}}
  {\hrulefill}


% --- Shortcuts ---
\newcommand{\R}{\mathbb{R}}
\newcommand{\inner}[2]{\left\langle #1,#2 \right\rangle}
\newcommand{\proj}{\mathrm{proj}}
\newcommand{\Id}{\mathbf{I}}
\newcommand{\ot}{\otimes}
\newcommand{\mc}{\mathcal}

\begin{document}
\noindent CS 58500 Theoretical Computer Science Toolkit (Spring 2026)\hfill Problem Set \# 2\\
%\rightline{\textbf{Due: September 15th}}
\{\{Your name\}\} \hfill \textbf{Due: Mar 2nd}

\noindent\hrulefill

The questions have been labeled with the date of the lecture in which the relevant material is covered.

\subsection*{Problem 1 (02/12, 02/17). (30 pts)}
\begin{enumerate}
\item For a matrix $A$, vector $b$, positive scalars $t_1, t_2 \geq 0$, and convex functions $f_1$ and $f_2$, the
function $g(x) = t_1f_1(Ax + b) + t_2f_2(x)$ is convex.
\item Show that the intersection of convex sets is convex; show that for convex functions, $f, g$, the function $h(x) = \max\{f(x), g(x)\}$ is also convex.
\item Show that for any $t\geq 0$, the level set $L(t):= \{x : f(x) \geq t\}$ of a logconcave function $f$ is convex.
\end{enumerate}

\noindent{\bfseries Solution.} \newline
\   %%% <-- ERASE THIS LINE AND WRITE YOUR SOLUTION HERE
\newpage
\subsection*{Problem 2 (02/19). (20 pts)}
Prove the following Theorem of Bieberbach from 1915: Let $K\subseteq \mathbb{R}^n$ be a compact set. Then
\begin{align*}
    \mathrm{Vol}_n(K)\leq \mathrm{Vol}_n(B_2^n)\left(\frac{\mathrm{diam}(K)}{2}\right)^n
\end{align*}
where $\mathrm{diam}(K):=\max\{\|x-y\|_2:x,y\in K\}$.\\
\begin{enumerate}
    \item Show that $\mathrm{diam}(\mathrm{conv}(K))=\mathrm{diam}(K)$. So, we can assume that $K$ is convex.
    \item For any unit vector $u\in \mathbb{S}^{n-1}$, let $S_u(K)$ be the Steiner symmetral of $K$ along $u$. For each $x\in u^\bot:=\{y\in \mathbb{R}^n:\langle y,u\rangle=0\}$, the length of the segment $K\cap (x+\mathbb{R}u)$ equals the length of $S_u(K)\cap (x+\mathbb{R}u)$. Steiner symmetrization can preserve convexity and volume. In this sub-problem, you need to prove that Steiner symmetrization cannot increase the diameter:
    \begin{align*}
        \mathrm{diam}(S_u(K))\leq \mathrm{diam}(K)
    \end{align*}
    Hint. Consider $M=\frac{K+\rho K}{2}$, where $\rho$ is the reflection across the hyperplane $u^\bot$.
    \item Show that sequentially applying Steiner symmetrizing w.r.t. the standard basis ($u_1=e_2,u_2=e_2,\dots,u_n=e_n$) makes the set centrally symmetric (i.e., $K=-K$).
    \item Prove the volume bound in the Bieberbach Theorem.
\end{enumerate}
\noindent{\bfseries Solution.} \newline
\   %%% <-- ERASE THIS LINE AND WRITE YOUR SOLUTION HERE
\newpage
\subsection*{Problem 3 (02/12, 02/24). (30 pts)}
Let $f:\mathbb{R}^n\rightarrow \mathbb{R}$ be differentiable and convex. We say $f$ is $L$-smooth if
\begin{align*}
    |\nabla f(x) -\nabla f(y)|\leq L\|x-y\|_2~~~\forall x,y\in \mathbb{R}^n
\end{align*}
\begin{enumerate}
    \item Prove that $f$ is $L$-smooth if and only if
    \begin{align*}
        f(y)\leq f(x)+\langle \nabla f(x),y-x\rangle +\frac{L}{2}\|y-x\|_2^2~~~\forall x,y\in \mathbb{R}^n        
    \end{align*}
    \item If $f$ is twice-differentiable, $f$ is $L$-smooth if and only if 
    \begin{align*}
        |\nabla^2f(x)[v,v]|=|v^\top \nabla^2f(x) v|\leq L\|v\|_2^2~~~\forall x,v\in \mathbb{R}^n
    \end{align*}
\end{enumerate}

\noindent{\bfseries Solution.} \newline
\   %%% <-- ERASE THIS LINE AND WRITE YOUR SOLUTION HERE
\newpage
\subsection*{Problem 4. (20 pts)}
The spanning tree polytope of an undirected graph $G = (V,E)$ is given by
\begin{align*}
    P=\{x\in \mathbb{R}_{\geq 0}^{|E|}:\sum_{e\in E}x_e=|V|-1, \sum_{(u,v)\in E\cap (S\times S)}x_{u,v}\leq |S|-1~~~\forall S\subseteq V,|S|\geq 1\}
\end{align*}
Then, the extreme points of this polytope exactly correspond to the indicator vectors of spanning
trees of $G$. 

In this problem, you need to design a membership oracle for this polytope. That is, for any input vector $x\in \mathbb{R}^n$, decide whether $x\in P$ in $\mathrm{poly}(|V|,|E|)$ time.

\noindent{\bfseries Solution.} \newline
\   %%% <-- ERASE THIS LINE AND WRITE YOUR SOLUTION HERE
\newpage
\end{document}

